\documentclass[11pt,a4paper]{article}
\usepackage[provide=*,spanish]{babel}
\usepackage[utf8]{inputenc}
\usepackage[T1]{fontenc}
\usepackage{lmodern,amsmath,booktabs,tabularx,xcolor,listings,microtype,parskip,hyperref,titlesec}
\usepackage[margin=2.35cm]{geometry}
\setlength{\emergencystretch}{3em}
\definecolor{azulWCD}{RGB}{17,67,112}
\definecolor{tituloColor}{RGB}{0,55,110}
\definecolor{reglaColor}{RGB}{0,90,160}
\hypersetup{colorlinks=true,linkcolor=azulWCD,pdftitle={Propiedades ópticas de los medios del WCD},pdfauthor={Jaime A. Betancourt}}
\lstset{language=C++,basicstyle=\ttfamily\footnotesize,breaklines=true,frame=single,backgroundcolor=\color{gray!6}}
\newcommand{\codigo}[1]{\texttt{#1}}
\titleformat{\section}{\normalfont\Large\bfseries\color{azulWCD}}
  {\thesection.}{0.6em}{}
\titleformat{\subsection}{\normalfont\large\bfseries}
  {\thesubsection.}{0.5em}{}

\begin{document}
\pagestyle{empty}

%------------------------------------------------
% PORTADA
%------------------------------------------------

\begin{titlepage}
\begin{center}

\vspace*{0.5cm}
\textcolor{reglaColor}{\rule{\linewidth}{1.2pt}}

\vspace{1.3cm}
{\Large\scshape Universidad Industrial de Santander}\\[0.2cm]
{\large Escuela de Física --- Facultad de Ciencias}\\[0.2cm]
{\large Doctorado en Física}

\vspace{2.2cm}
{\Huge\bfseries\color{tituloColor} Propiedades ópticas\\[0.2cm]
de los medios del WCD}

\vspace{1.8cm}
{\large\textbf{Informe técnico}}

\vspace{0.6cm}
\begin{minipage}{0.82\linewidth}
\centering
\textit{``Implementación de las propiedades ópticas de los medios agua--NaCl en el WCD''}
\end{minipage}

\vfill

{\large\textbf{Autor}}\\[0.15cm]
Jaime A. Betancourt\\[0.1cm]
Doctorado en Física, Universidad Industrial de Santander

\vspace{1.2cm}
\textcolor{reglaColor}{\rule{0.5\linewidth}{0.6pt}}

\vspace{1cm}
11 de septiembre de 2026

\vspace{1cm}
\textcolor{reglaColor}{\rule{\linewidth}{1.2pt}}

\end{center}
\end{titlepage}

\pagestyle{plain}

\begin{abstract}
Se documenta la ampliación de \codigo{Materials-mejorado.cc} para asignar una
tabla de propiedades ópticas independiente al agua pura y a las soluciones con
2.5, 5.0 y 10.0\,\% de NaCl. Cada tabla combina un índice de refracción
espectral específico con la longitud de absorción suministrada y conserva,
como aproximación explícita, el modelo Mie preexistente del ejemplo OpNovice de
Geant4. \codigo{SaltyWCD-mejorado.cc} selecciona el material completo y
comprueba que la tabla óptica esté asociada al volumen activo.
\end{abstract}

\section{Malla espectral}
Las cuatro matrices de índice y la longitud de absorción contienen 32 valores.
Se asociaron a la matriz creciente \codigo{water2PhotonEnergy}, también de 32
elementos, que cubre de 2.034 a 4.136 eV. La correspondencia entre energía y
longitud de onda se obtiene con
\begin{equation}
 E_\gamma[\mathrm{eV}]=\frac{1239.841984}{\lambda[\mathrm{nm}]}.
\end{equation}
No se empleó \codigo{water1PhotonEnergy}, porque tiene 30 elementos y habría
producido una asociación dimensional y espectral incorrecta.

\section{Índice de refracción}
Se incorporaron cuatro dispersiones, almacenadas respectivamente en
\codigo{pureWaterRefIndex}, \codigo{saltyWater2p5RefIndex},
\codigo{saltyWater5RefIndex} y
\codigo{saltyWater10RefIndex}.
La interpolación lineal que realiza la tabla de propiedades de Geant4 entre los
puntos vecinos de 2.103 y 2.139 eV produce:
\begin{center}
\begin{tabular}{lc}
\toprule
Medio & $n(589.3\,\mathrm{nm})$\\
\midrule
Agua pura & 1.3330000\\
Agua + 2.5\,\% NaCl & 1.3396957\\
Agua + 5.0\,\% NaCl & 1.3435863\\
Agua + 10.0\,\% NaCl & 1.3593570\\
\bottomrule
\end{tabular}
\end{center}
Los resultados reproducen los valores objetivo indicados.

\section{Longitud de absorción}
La matriz \codigo{waterAbsLen} se añadió sobre la misma malla de 32 energías.
En ausencia de matrices de absorción medidas por separado para cada
concentración, se usa como línea base común en los cuatro medios. Esta decisión
no afirma que la salinidad carezca de efecto sobre la absorción; evita inventar
una dependencia espectral no suministrada y mantiene visible la aproximación.

\section{Dispersión óptica}
Se conserva \codigo{MIEHG} con \codigo{water2PhotonEnergyMie},
\codigo{water2Mie} y las constantes direccionales preexistentes. El mismo
modelo se conecta provisionalmente a las cuatro tablas. Corresponde al ejemplo
OpNovice y no constituye una medición específica de las soluciones salinas.
Cuando se disponga de datos de dispersión en función de la concentración, esta
propiedad deberá reemplazarse por tablas independientes.

No se añadió manualmente \codigo{RAYLEIGH}: Geant4 puede calcular la
dispersión Rayleigh para agua cuando dispone de las propiedades termodinámicas
requeridas, pero extender ese cálculo a soluciones salinas exige parámetros
validados. Añadir una longitud arbitraria introduciría una hipótesis adicional.

\section{Construcción de las tablas}
Se creó una única función auxiliar:
\begin{lstlisting}
static G4MaterialPropertiesTable* CreateWCDOpticalTable(
    G4double* refractiveIndex);
\end{lstlisting}
Cada llamada crea un objeto distinto y añade \codigo{RINDEX},
\codigo{ABSLENGTH}, \codigo{MIEHG}, \codigo{MIEHG\_FORWARD},
\codigo{MIEHG\_BACKWARD} y \codigo{MIEHG\_FORWARD\_RATIO}. El agua recibe
\codigo{waterPT1}; cada solución recibe su propia tabla mediante
\codigo{CreateSaltyWaterByFinalMassFraction}. Por tanto, no se comparte el
puntero de la tabla entre concentraciones.

\codigo{waterPT2} se conserva como alias de \codigo{waterPT1} para evitar
romper código externo que todavía consulte ese miembro. Esta compatibilidad no
crea una segunda definición física.

\section{Conexión con SaltyWCD.cc}
\codigo{SelectWCDMedium} continúa escogiendo el medio por fracción másica
final. Como la tabla se asigna al construir cada \codigo{G4Material}, el
\codigo{G4LogicalVolume} del tanque recibe simultáneamente composición,
densidad e información óptica coherentes.

Se añadió un diagnóstico que consulta
\codigo{GetMaterialPropertiesTable()} e informa si las propiedades
\codigo{RINDEX/ABSLENGTH/MIEHG} están asignadas. No se duplican matrices
ópticas en la clase geométrica.

\section{Propiedades que permanecen independientes}
No se modificaron:
\begin{itemize}
\item la reflectividad, retrodispersión y componentes especulares del Tyvek;
\item el índice y la absorción del Pyrex;
\item la eficiencia cuántica o la respuesta del PMT;
\item las densidades y fracciones másicas implementadas previamente;
\item las listas y procesos físicos de Geant4.
\end{itemize}
Estas propiedades pertenecen a superficies, ventanas o dispositivos distintos
del medio activo.

\section{Comprobaciones realizadas y pendientes}
La revisión estática confirmó 32 elementos en cada matriz, energía estrictamente
creciente, longitudes de absorción positivas y los cuatro valores interpolados
en 589.3 nm. También se ejecutó \codigo{git diff --check}.

Antes de producción se requiere compilar con MEIGA y Geant4 10.07.p04 y ejecutar
corridas cortas para los cuatro medios. Deben verificarse la creación de
fotones Cherenkov, su transporte y llegada al fotocátodo, además de revisar que
el mensaje de diagnóstico indique \codigo{assigned}. La comparación física
debe mantener constantes geometría, fuente, lista de física, número de
incidentes y respuesta del PMT.

\section{Archivos propuestos}
\begin{itemize}
\item \codigo{Materials-mejorado-optica.cc};
\item \codigo{SaltyWCD-mejorado-optica.cc}.
\end{itemize}
Para integrarlos deben renombrarse como \codigo{Materials.cc} y
\codigo{SaltyWCD.cc} en la ruta usada por \codigo{G4WCDSimulator}, después de
revisar el diff contra la rama activa.
\end{document}
